\documentclass[12pt,a4paper]{article}
\usepackage[utf8]{inputenc}
\usepackage[czech]{babel}
\usepackage[T1]{fontenc}
\usepackage{amsmath}
\usepackage{amsfonts}
\usepackage{amssymb}
\usepackage{graphicx}
\usepackage{titlesec}
\usepackage[left=2cm,right=2cm,top=2cm,bottom=2cm]{geometry}
\usepackage{indentfirst}
\usepackage{listings}
\usepackage{color}
\usepackage{multirow}

%Pravidlo pro řádkování
\renewcommand{\baselinestretch}{1.5}

%Pravidlo pro začínání kapitol na novém řádku
\let\oldsection\section
\renewcommand\section{\clearpage\oldsection}

%Formáty písem pro nadpisy (-změněno na bezpatkové \sffamily z původního \normalfont
\titleformat{\section}
{\sffamily\Large\bfseries}{\thesection}{1em}{}
\titleformat{\subsection}
{\sffamily\large\bfseries}{\thesubsection}{1em}{}
\titleformat{\subsubsection}
{\sffamily\normalsize\bfseries}{\thesubsubsection}{1em}{}

%Nastavení zvýrazňování kódu v \lslisting
\definecolor{mygreen}{rgb}{0,0.6,0}
\definecolor{mygray}{rgb}{0.5,0.5,0.5}
\lstset{commentstyle=\color{mygreen},keywordstyle=\color{blue},numberstyle=\tiny\color{mygray}}

\author{Jan Šmejkal}

\begin{document}

%-------------Úvodni strana---------------
\begin{titlepage}

\includegraphics[width=50mm]{img/FAV.jpg}
\\[160 pt]
\centerline{ \Huge \sc KIV/ZOS - Základy operačních systémů}
\centerline{ \huge \sc Semestrální práce }
\\[12 pt]
{\large \sc
\centerline{Zadání pro login kiwi}
}


{
\vfill 
\parindent=0cm
\textbf{Jméno:} Štěpán Ševčík\\
\textbf{Osobní číslo:} A13B0443P\\
\textbf{E-mail:} kiwi@students.zcu.cz\\
\textbf{Datum:} {\large \today\par} %datum

}

\end{titlepage}

%------------------------------------------

%------------------Obsah-------------------
\newpage
\setcounter{page}{2}
\setcounter{tocdepth}{3}
\tableofcontents
%------------------------------------------

%--------------Text dokumentu--------------
\section{Zadání}
\subsection{1. část zadání}
Kontrola, zda každý řetěz FAT má správnou délku (odpovídá velikosti souboru v adresáři)
\subsection{2. část zadání}
Setřesení volného místa (volné místo na konci, ale bloky souborů nemusí jít za sebou).
\subsection{Další požadavky na aplikaci}
\begin{description}
\item Program vytvoří více vláken s možností zadat při/po spuštění počet paralelních entit
\item Zdrojový kód programu bude komentovaný - např. co dělá každá funkce
\item Bude ošetřeno špatné zadání vstupů
\item Součástí řešení bude skript v bashi, který umožní automatické spuštění všech testů nutných k vytvoření dokumentace (tj. např. spuštění s různým počtem vláken)
\item Dokumentace v elektronické podobě, popisující zejména použité konstrukce pro vytváření procesů/vláken/synchronizaci a použité algoritmy
\item
\end{description}

\section{Rozbor řešení úlohy}
\subsection{Spouštěč úlohy}
Aplikace obsahuje čtyři rutiny, mezi kterými lze vybírat požadovanou pomocí spouštěcích argumentů. Dvě rutiny jsou optimalizované verze zadaných rutin \textit{fat\_reader} a \textit{fat\_writer} a zbývající dvě rutiny vykonávají zadané úkony. Při spuštění se pomocí dedikované funkce určí požadovaná procedura a případně i počet vláken. Pokud je zadána neplatná kombinace parametrů nebo pokud je první parametr \textbf{help}, na obrazovku se vypíše instrukce pro použití a výkon programu skončí. Jinak se vykoná požadovaná instrukce definovaná v odpovídajícím souboru.
\subsection{1. část: Kontrola délky řetězce FAT}
Ověření správné délky řetězce fat lze provést dvěma různými způsoby: hrubě podle počtu clusterů, kde se uložená délka řetězce vydělí délkou clusteru a výsledek se porovná s počtem clusterů napočítaných v řetězu od začátku souboru do značky FAT\_FILE\_END, anebo jemně podle znaků v jednotlivých clusterech celého řetězce. Svoje řešení jsem vypracovával na jemném základu a celkovou délku souboru vypočítávám s přesností na znaky. Pokud vycházím z předpokladu že se clustery vyplňují od začátku a nenačíná se další cluster, pokud předchozí není zaplněný, stačí spočítat clustery, které nejsou v řetězci konečné, vynásobit jejich počet stanovenou délkou clusteru a připočíst k mezivýsledku délku posledního clusteru řetězce zjištěnou pomocí funkce \texttt{strlen}.
Při kontrole délky se také kontroluje možné zacyklení řetězce od momentu, kdy se při procházení clusterů zjistí větší počet než očekávaný počet clusterů.
Bezpečnost paralelního výpočtu zajišťuje návrh ve stylu farmer worker, kdy se takzvaný farmer stará o přidělování souborů a jednotliví workeři provádějí výpočty. V tomto návrhovém vzoru nastávají dvě kritické sekce při čtení ze souboru: jedna je načítání \texttt{root\_directory} záznamů a druhá nastává při načítání clusterů řetězce. Obě tyto kritické sekce jsou řešeny mutexem, který hlídá přístup do souboru.
\newpage
\subsection{2. část: Setřesení volného místa}
Setřásání volného místa funguje opět podle návrhového vzoru farmer-worker, kdy jednotlivá vlákna náležící workerům žádají farmáře o přidělení práce. Farmář se o přidělení práce rozhoduje na základě předem provedené analýzy tabulky FAT, kdy si postupně prohlédl všechny neprázdné záznamy o clusterech a rozdělí si je do shluků (tzv. cluster chunků) jednotné velikosti. Každý worker poté prochází svůj interval a hledá clustery, které postupně skládá ve svém chunku, který je jednoznačně dán jeho pořadovým indexem. Workeři také dále také kontrolují, zda-li v cílovém clusteru ještě není volno a pokud se tam vyskytuje nepřesunutý cluster, čekají dokud se místo neuvolní.
Kritických sekcí se při přesunu vyskytuje hned několik, protože je potřeba zachovávat řetěz a přesouvaný cluster spolu také musí zamknout předcházející a následující clustery v řetězci, pokud se nejedná o první či poslední článek řetězu.
S přesuny jednotlivých clusterů také vzniká riziko ztráty reference začátku řetězce. Toto riziko je ošetřeno zpětnou tabulkou indexů, která popisuje každý cluster
Po přesunutí všech clusterů se uloží výsledná FAT a root directory záznamy.

\section{Měření výkonu}
\begin{tabular}{ l r }

\begin{tabular}{| c || c | c |}
\hline
\multirow{2}{*}{Pož. vláken} & \multicolumn{2}{c |}{Délka vykonávání} \\
				& Kontrola délky & Setřesení \\ \hline \hline
1 & 0.52 ms & 3.49 ms \\ \hline
2 & 0.46 ms & 3.85 ms \\ \hline
3 & 0.42 ms & 3.84 ms \\ \hline 
4 & 0.53 ms & 3.69 ms \\ \hline
5 & 0.43 ms & 4.06 ms \\ \hline
6 & 0.47 ms & 4.31 ms \\ \hline 
7 & 0.55 ms & 4.32 ms \\ \hline
8 & 0.58 ms & 4.30 ms \\ \hline	

\end{tabular}
&
\begin{tabular}{| c | c |}
\hline
\multicolumn{2}{|c |}{Konfigurace VB stroje} \\ \hline \hline
OS & Debian 64-bit \\ \hline
Operační paměť & 1024 MB \\ \hline
Přiřazených CPU & 2 \\ \hline
\multicolumn{2}{c}{} \\
\hline
\multicolumn{2}{| c |}{Systémové prostředky} \\ \hline \hline
CPU & Intel Core i7 \\ \hline
Frekvence CPU & 2.40GHz \\ \hline
Počet jader & 2 \\ \hline
Paměť & 1024 MB \\ \hline

\end{tabular}
\end{tabular}
\section{Závěr}
Program jsem vyvíjel strukturovaně a snažil jsem se o vysokou úroveň modulárnosti. Vytváření algoritmu a sémantiky zamykání trvalo přibližně stejně dlouho jako jeho implementace ale díky tomuto jsem se po většinou nesetkal s 'nesmyslnými' výsledky. Vypracovaná část zadání zcela podporuje paralelní vykonání ale na dostupných testovacích datech se rozdílný počet vláken výrazně neprojevuje.


%------------------------------------------

\end{document}